\documentclass{article}

\usepackage[a4paper,margin=2cm]{geometry}
\usepackage{../../../../components/mathtex}
\usepackage{colortbl}
\usepackage{array}
\usepackage{hyperref}
\usepackage{stmaryrd}

\usepackage{fancyhdr}
\pagestyle{fancy}
\fancyhf{}
\fancyhead[L]{DL3 Math-Info}
\fancyhead[C]{Révisions de la Toussaint}
\fancyhead[R]{2026-2027}
\fancyfoot[L]{Ewen Rodrigues de Oliveira}
\fancyfoot[R]{\thepage}

\begin{document}

\thispagestyle{empty}

\newcommand{\doctitle}{Révisions pour les vacances de la Toussaint}

\vspace*{\fill}

\begin{center}
    {\LARGE\textbf{\doctitle}}\\[0.6em]
    {\Large Double licence 3 Mathématiques et Informatique  Fondamentale}\\[2em]
    {\large Année universitaire 2026-2027}\\[0.25em]
    {\large par \textit{Ewen Rodrigues de Oliveira}}
\end{center}

\vfill

\noindent
Ce livret propose des révisions sur les matières au programme du semestre 5 de la double licence. Je n'ai pas créé ces exercices : j'ai surtout rassemblé des exercices intéressants issus d'annales et de feuilles de TD de l'\href{https://u-pariscite.fr/}{Université Paris Cité}.

\medskip

\noindent
Le programme n'est sans doute pas exhaustif. S'il y a une erreur, une coquille ou un énoncé ambigu, merci de me contacter pour que je puisse corriger cela rapidement.

\medskip

\noindent
Si besoin d'un corrigé, réferrez-vous à \href{https://ewenrdo.fr/ressources}{mon site internet}. Réciproquement, si vous avez un corrigé à proposer, n'hésitez pas à me le faire parvenir.

\remark{Je conseille de faire les exercices avec son cours à proximité. Il peut être nécessaire de passer par des résultats vus en cours pour résoudre les exercices, et c'est aussi l'occasion de les revoir.}

\medskip
\noindent
\textbf{Le livret peut être utilisé par les étudiants en licence simple (mathématiques resp. informatique), il suffit de ne pas faire les exercices qui ne sont pas dans le programme de la licence simple.}

\begin{flushright}
\textbf{Bonnes vacances \faSmile}
\end{flushright}

\newpage

\section*{Crédits}
\begin{itemize}
\item \textit{Non rédigé pour le moment.}
\end{itemize}

\tableofcontents

\newpage
\renewcommand{\arraystretch}{1.3}
\section{Jour 1}
\newpage
\section{Jour 2}
\newpage
\section{Jour 3}
\newpage
\section{Jour 4}
\newpage
\section{Jour 5}
\newpage
\section{Jour 6}
\newpage
\section{Jour 7}
\newpage
\section{Jour 8}
\newpage
\section{Jour 9}
\newpage

\end{document}
